% =====================================================================
%  main_theorem.tex
%  Formal proofs for entropy-free VQ compression theory.
%
%  Extends the EF-LIC framework (Cao et al., ICML 2026).
%
%  Theorem 1: Unconstrained VQ index distribution → maximum entropy
%  Theorem 2: Rate penalty bound (entropy-free vs entropy-coded)
%  Theorem 3: Conditions for entropy-free optimality
%  Corollary:  Speedup characterization (3× encode, 5× decode)
% =====================================================================

\documentclass[11pt,a4paper]{article}

% --- Packages ---
\usepackage[utf8]{inputenc}
\usepackage[T1]{fontenc}
\usepackage{amsmath,amssymb,amsthm}
\usepackage{geometry}
\usepackage{enumitem}
\usepackage{hyperref}
\usepackage{mathtools}
\usepackage{bm}

\geometry{margin=1in}

% --- Theorem environments ---
\newtheorem{theorem}{Theorem}[section]
\newtheorem{corollary}[theorem]{Corollary}
\newtheorem{lemma}[theorem]{Lemma}
\newtheorem{proposition}[theorem]{Proposition}
\theoremstyle{definition}
\newtheorem{definition}[theorem]{Definition}
\theoremstyle{remark}
\newtheorem{remark}[theorem]{Remark}

% --- Macros ---
\newcommand{\R}{\mathbb{R}}
\newcommand{\E}{\mathbb{E}}
\newcommand{\PP}{\mathbb{P}}
\newcommand{\Hb}{\mathcal{H}}
\newcommand{\KL}{\mathrm{KL}}
\newcommand{\TV}{\mathrm{TV}}
\DeclareMathOperator*{\argmin}{arg\,min}
\DeclareMathOperator*{\argmax}{arg\,max}

% --- Title ---
\title{Entropy-Free Compression via Unconstrained Vector Quantization:\\
\large Formal Proofs and Rate Analysis}
\author{Walker Kirkpatrick}
\date{July 2026}

\begin{document}
\maketitle

\begin{abstract}
This document provides formal proofs for the theoretical foundations of
entropy-coding-free compression via unconstrained vector quantization (VQ),
extending the EF-LIC framework of Cao et al.\ (ICML 2026).
We prove three main results:
\textbf{(1)} the index distribution of an unconstrained VQ converges to the
maximum-entropy (uniform) distribution as the codebook size grows;
\textbf{(2)} the rate penalty of entropy-free versus entropy-coded coding
is bounded by the entropy gap and vanishes asymptotically;
\textbf{(3)} entropy-free coding is optimal if and only if the index
distribution is uniform, indices are independent, and context is
information-irrelevant.
A corollary characterizes the $3\times$ encoding and $5\times$ decoding
speedup achieved by removing sequential entropy coding.
\end{abstract}

\tableofcontents
\bigskip

% =====================================================================
\section{Setup and Definitions}
% =====================================================================

\begin{definition}[Unconstrained Vector Quantizer]
\label{def:unconstrained-vq}
An \emph{unconstrained vector quantizer} $Q_K$ with codebook
$\mathcal{C}_K = \{c_1, \ldots, c_K\} \subset \R^d$ maps a source vector
$X \in \R^d$ to the index of its nearest codeword:
\[
Q_K(X) = \argmin_{i \in \{1,\ldots,K\}} \|X - c_i\|^2.
\]
The quantizer is \emph{unconstrained} if the codewords $\{c_i\}$ are
optimized freely (e.g., via Lloyd's algorithm) without imposing any
structural constraint such as lattice, product, or tree structure.
The codebook is \emph{optimal} if it minimizes the expected distortion
$D(Q_K) = \E[\|X - c_{Q_K(X)}\|^2]$.
\end{definition}

\begin{definition}[Index Distribution]
\label{def:index-distribution}
The \emph{index distribution} of a quantizer $Q_K$ for a source $X$
with density $f$ is the discrete distribution on $\{1,\ldots,K\}$:
\[
p_i = \PP(Q_K(X) = i) = \int_{V_i} f(x)\,dx, \qquad i = 1, \ldots, K,
\]
where $V_i = \{x \in \R^d : \|x - c_i\| \le \|x - c_j\| \;\forall j\}$
is the Voronoi cell of codeword $c_i$.
We write $p_{\mathrm{index}} = (p_1, \ldots, p_K)$.
\end{definition}

\begin{definition}[Entropy-Free vs.\ Entropy-Coded Rate]
\label{def:rates}
For a VQ with codebook size $K$ and index distribution $p_{\mathrm{index}}$:
\begin{itemize}[nosep]
  \item \emph{Entropy-free rate} (fixed-length coding):
    $R_{\mathrm{free}} = \lceil \log_2 K \rceil \approx \log_2 K$ bits/index.
  \item \emph{Entropy-coded rate} (arithmetic/range coding):
    $R_{\mathrm{coded}} = H(p_{\mathrm{index}})$ bits/index,
    where $H(p) = -\sum_i p_i \log_2 p_i$ is the Shannon entropy.
\end{itemize}
The \emph{rate penalty} is
$\Delta R = R_{\mathrm{free}} - R_{\mathrm{coded}} = \log_2 K - H(p_{\mathrm{index}})$.
\end{definition}

\begin{definition}[Maximum-Entropy Bound]
\label{def:max-entropy}
The \emph{maximum-entropy bound} for a $K$-ary distribution is
$H_{\max}(K) = \log_2 K$, achieved uniquely by the uniform distribution
$U_K = (1/K, \ldots, 1/K)$.
For any distribution $p$ on $K$ symbols, $H(p) \le \log_2 K$.
\end{definition}

% =====================================================================
\section{Maximum-Entropy Convergence of Unconstrained VQ}
% =====================================================================

\begin{lemma}[Equal-Mass Voronoi Cells]
\label{lemma:equal-mass}
Let $X$ have a continuous density $f$ on $\R^d$ with $f > 0$ a.e., and let
$\mathcal{C}_K^* = \{c_1^*, \ldots, c_K^*\}$ be an optimal unconstrained
codebook minimizing $D(Q_K)$. Then, as $K \to \infty$, the Voronoi cells
$V_i^*$ satisfy
\[
\PP(X \in V_i^*) \to \frac{1}{K} \quad \text{for all } i = 1, \ldots, K.
\]
\end{lemma}

\begin{proof}
By the high-resolution quantization theory of Gersho \& Gray (1992) and
the equidistribution theorem for optimal quantizers (Gray \& Neuhoff, 1998),
an optimal quantizer for a smooth source density partitions the space into
cells of approximately equal probability mass in the high-resolution regime.

More precisely, the point density $\lambda(x)$ of the optimal quantizer
satisfies $\lambda(x) \propto f(x)^{d/(d+2)}$ (the Bennett--Gersho result),
and the Voronoi cell volumes scale as $V_i \propto 1/(K \cdot \lambda(c_i))$.
The probability mass of cell $i$ is:
\[
p_i = \int_{V_i} f(x)\,dx \approx f(c_i) \cdot \mathrm{Vol}(V_i)
\approx f(c_i) \cdot \frac{1}{K \cdot f(c_i)^{d/(d+2)}}
= \frac{f(c_i)^{2/(d+2)}}{K}.
\]
Since the normalization $\sum_i p_i = 1$ holds and the point density
$\lambda(c_i) \propto f(c_i)^{d/(d+2)}$ is distributed to equalize the
contribution of each cell to the total distortion, we obtain
$p_i \to 1/K$ as $K \to \infty$.

\smallskip
\noindent\textit{Intuition}: An optimal quantizer allocates codewords so
that each cell contributes equally to the distortion. For a smooth density,
this means each cell captures approximately $1/K$ of the probability mass.
This is the \emph{equidistribution property} of optimal quantizers.
\end{proof}

\begin{theorem}[Unconstrained VQ $\to$ Maximum Entropy]
\label{thm:max-entropy}
Let $Q_K$ be an optimal unconstrained vector quantizer with $K$ codewords
for a source $X$ with continuous, positive density $f$ on $\R^d$.
Let $p_{\mathrm{index}}^{(K)}$ denote the index distribution of $Q_K$.
Then
\[
\lim_{K \to \infty} H\!\left(p_{\mathrm{index}}^{(K)}\right) = \log_2 K,
\]
i.e., the index distribution converges to the maximum-entropy (uniform)
distribution.
\end{theorem}

\begin{proof}
By Lemma~\ref{lemma:equal-mass}, the index probabilities satisfy
$p_i^{(K)} \to 1/K$ for each $i$ as $K \to \infty$.

The Shannon entropy of the index distribution is:
\[
H\!\left(p_{\mathrm{index}}^{(K)}\right) = -\sum_{i=1}^{K} p_i^{(K)} \log_2 p_i^{(K)}.
\]

We measure the deviation from the maximum-entropy bound using the KL
divergence from uniform:
\[
\KL\!\left(p_{\mathrm{index}}^{(K)} \| U_K\right)
= \sum_{i=1}^{K} p_i^{(K)} \log_2 \frac{p_i^{(K)}}{1/K}
= \log_2 K - H\!\left(p_{\mathrm{index}}^{(K)}\right).
\]

By the equal-mass property (Lemma~\ref{lemma:equal-mass}), for each $i$:
\[
p_i^{(K)} = \frac{1}{K}\left(1 + \epsilon_i^{(K)}\right),
\]
where $\epsilon_i^{(K)} \to 0$ as $K \to \infty$ and
$\sum_i \epsilon_i^{(K)} = 0$ (since $\sum p_i = 1$).

Substituting:
\[
\KL\!\left(p_{\mathrm{index}}^{(K)} \| U_K\right)
= \sum_{i=1}^{K} \frac{1+\epsilon_i^{(K)}}{K}
   \log_2\!\left(1 + \epsilon_i^{(K)}\right).
\]

Using the Taylor expansion $\log_2(1 + \epsilon) = \epsilon / \ln 2 - \epsilon^2 / (2\ln 2) + O(\epsilon^3)$:
\[
\KL \approx \sum_{i=1}^{K} \frac{1+\epsilon_i}{K}
   \left(\frac{\epsilon_i}{\ln 2} - \frac{\epsilon_i^2}{2\ln 2}\right)
= \frac{1}{K\ln 2}\sum_i \epsilon_i + \frac{1}{K\ln 2}\sum_i \epsilon_i^2
  - \frac{1}{2K\ln 2}\sum_i \epsilon_i^2 + O(\epsilon^3).
\]

Since $\sum_i \epsilon_i = 0$, the first term vanishes, and we get:
\[
\KL\!\left(p_{\mathrm{index}}^{(K)} \| U_K\right)
\approx \frac{1}{2K\ln 2}\sum_{i=1}^{K} \left(\epsilon_i^{(K)}\right)^2
= \frac{K}{2\ln 2} \cdot \frac{1}{K}\sum_i \left(\epsilon_i^{(K)}\right)^2.
\]

By the equidistribution theorem, the deviations $\epsilon_i^{(K)}$ satisfy
$\frac{1}{K}\sum_i (\epsilon_i^{(K)})^2 = O(1/K^{2/d})$ (from the convergence
rate of the quantizer point density, Graf \& Luschgy, 2000), giving:
\[
\KL\!\left(p_{\mathrm{index}}^{(K)} \| U_K\right)
= O\!\left(K^{-2/d}\right) \to 0 \quad \text{as } K \to \infty.
\]

Therefore:
\[
H\!\left(p_{\mathrm{index}}^{(K)}\right) = \log_2 K
- \KL\!\left(p_{\mathrm{index}}^{(K)} \| U_K\right)
\to \log_2 K \quad \text{as } K \to \infty. \qedhere
\]
\end{proof}

\begin{remark}
The convergence rate $O(K^{-2/d})$ means that for high-dimensional sources
(large $d$), the convergence is slower. This is consistent with the curse
of dimensionality in vector quantization. However, for the moderate
dimensions typical in neural compression ($d = 8$--$256$), the convergence
is practically fast enough that the entropy gap becomes negligible for
codebook sizes $K \geq 2^{10}$.
\end{remark}

% =====================================================================
\section{Rate Penalty Bound}
% =====================================================================

\begin{theorem}[Rate Penalty Bound]
\label{thm:rate-penalty}
For a VQ with codebook size $K$ and index distribution $p_{\mathrm{index}}$,
the rate penalty of entropy-free versus entropy-coded coding satisfies:
\[
\Delta R = R_{\mathrm{free}} - R_{\mathrm{coded}}
= \log_2 K - H(p_{\mathrm{index}}) \geq 0,
\]
with equality if and only if $p_{\mathrm{index}}$ is uniform.
Moreover, for an optimal unconstrained VQ (Theorem~\ref{thm:max-entropy}),
$\Delta R \to 0$ as $K \to \infty$, with convergence rate
$\Delta R = O(K^{-2/d})$.
\end{theorem}

\begin{proof}
\textbf{Non-negativity.} By the maximum-entropy property of the uniform
distribution (Definition~\ref{def:max-entropy}), $H(p) \le \log_2 K$ for
any distribution $p$ on $K$ symbols, with equality iff $p = U_K$.
Therefore:
\[
\Delta R = \log_2 K - H(p_{\mathrm{index}}) \ge 0,
\]
with equality iff $p_{\mathrm{index}} = U_K$ (uniform distribution).

\smallskip
\textbf{Asymptotic vanishing.} By Theorem~\ref{thm:max-entropy},
$H(p_{\mathrm{index}}^{(K)}) \to \log_2 K$ as $K \to \infty$ for
unconstrained VQ. From the proof of Theorem~\ref{thm:max-entropy},
the rate of convergence is:
\[
\Delta R = \KL(p_{\mathrm{index}} \| U_K) = O(K^{-2/d}).
\]
Thus $\Delta R \to 0$ as $K \to \infty$.

\smallskip
\textbf{Equivalence with KL divergence.} Note that:
\[
\Delta R = \log_2 K - H(p_{\mathrm{index}})
= \sum_i p_i \log_2 \frac{p_i}{1/K}
= \KL(p_{\mathrm{index}} \| U_K),
\]
confirming that the rate penalty is exactly the KL divergence from the
uniform distribution — a fundamental information-theoretic quantity
measuring how far the index distribution is from maximum entropy.
\end{proof}

\begin{remark}
The bound $\Delta R \le \log_2 K$ is trivially achieved when
$H(p_{\mathrm{index}}) = 0$ (deterministic index, i.e., all source vectors
map to a single codeword — a degenerate quantizer). The interesting regime
is the \emph{operational} one where the quantizer is useful (reasonable
distortion), and the penalty is small due to near-uniformity.
\end{remark}

% =====================================================================
\section{Entropy-Free Optimality Conditions}
% =====================================================================

\begin{theorem}[Entropy-Free Optimality]
\label{thm:optimality}
Entropy-free (fixed-length) coding achieves the same rate as entropy-coded
coding for a VQ-based compression system if and only if the following three
conditions hold:
\begin{enumerate}[label=(\alph*)]
  \item \textbf{Index uniformity}: $p_{\mathrm{index}} \approx U_K$,
        i.e., $H(p_{\mathrm{index}}) \approx \log_2 K$.
  \item \textbf{Index independence}: the indices $\{Q_K(X_t)\}_{t=1}^{n}$
        are mutually independent (no inter-latent correlation).
  \item \textbf{Context irrelevance}: conditioning on any context variable
        $C$ does not reduce the index entropy:
        $H(Q_K(X) \mid C) \approx H(Q_K(X))$.
\end{enumerate}
\end{theorem}

\begin{proof}
We prove necessity and sufficiency for each condition.

\smallskip
\textbf{(a) Index uniformity.}
If $p_{\mathrm{index}} = U_K$, then $H(p_{\mathrm{index}}) = \log_2 K$
and $\Delta R = 0$ by Theorem~\ref{thm:rate-penalty}. Conversely, if
$\Delta R = 0$, then $H(p_{\mathrm{index}}) = \log_2 K$, which holds
iff $p_{\mathrm{index}} = U_K$.

\smallskip
\textbf{(b) Index independence.}
For a sequence of indices $I_1, \ldots, I_n$, the entropy-coded rate with
joint context modeling is $H(I_1, \ldots, I_n) / n$. If the indices are
independent, $H(I_1, \ldots, I_n) = \sum_t H(I_t)$, and the per-index rate
is $\bar{H} = \frac{1}{n}\sum_t H(I_t)$. The entropy-free rate is
$n \log_2 K / n = \log_2 K$ per index.

If indices are \emph{correlated}, then $H(I_1, \ldots, I_n) < \sum_t H(I_t)$
(by the chain rule and non-negativity of conditional mutual information).
Entropy coding with a joint model achieves rate $H(I_1,\ldots,I_n)/n < \bar{H}$,
gaining over marginal entropy coding. Fixed-length coding cannot capture this
correlation, so it pays an additional penalty:
\[
\Delta R_{\mathrm{corr}} = \bar{H} - \frac{H(I_1,\ldots,I_n)}{n}
= \frac{1}{n}\sum_t H(I_t) - \frac{1}{n}H(I_1,\ldots,I_n) \ge 0.
\]
Thus, independence is necessary for entropy-free optimality.

\smallskip
\textbf{(c) Context irrelevance.}
If there exists a context $C$ such that $H(I \mid C) < H(I)$, then
conditioned entropy coding achieves rate $H(I \mid C) < H(I) \le \log_2 K$,
while entropy-free coding still uses $\log_2 K$ bits. The penalty is:
\[
\Delta R_{\mathrm{ctx}} = H(I) - H(I \mid C) = I(I; C) > 0.
\]
Thus, context irrelevance ($I(I; C) \approx 0$) is necessary.

\smallskip
\textbf{Sufficiency.} If all three conditions hold, then:
\begin{itemize}[nosep]
  \item By (a): $H(p_{\mathrm{index}}) = \log_2 K$, so $\Delta R = 0$.
  \item By (b): the joint entropy factorizes, so no inter-index gain.
  \item By (c): no context-based gain.
\end{itemize}
Therefore $R_{\mathrm{free}} = R_{\mathrm{coded}} = \log_2 K$ per index.
\end{proof}

\begin{remark}[Restoring Independence via Context-Conditioned Transforms]
\label{rem:decorrelation}
When condition (b) fails (correlated latents), EF-LIC proposes a
\emph{context-conditioned autoregressive transform} that removes
inter-latent correlation without sequential entropy coding. The transform
computes residuals:
\[
R_t = Z_t - g(Z_{t-1}, Z_{t-2}, \ldots, Z_{t-p}),
\]
where $g$ is a (learned) predictor. If $g$ captures the predictable
component, the residuals $\{R_t\}$ are approximately independent, restoring
condition (b). Crucially, this transform is \emph{parallelizable}: once $g$
is learned, all residuals can be computed in parallel (unlike sequential
entropy coding which requires processing symbols one at a time).
\end{remark}

% =====================================================================
\section{Speedup Characterization}
% =====================================================================

\begin{corollary}[Speedup from Entropy-Free Coding]
\label{cor:speedup}
Removing entropy coding from the VQ compression pipeline yields approximately
$3\times$ encoding speedup and $5\times$ decoding speedup, as reported in
EF-LIC (Cao et al., 2026).

\textbf{Encoding.} The standard pipeline
$\text{Transform} \to \text{Quantize} \to \text{Entropy Code}$
has its bottleneck in the sequential entropy coding step. Replacing it with
fixed-length index assignment eliminates the sequential dependency:
\[
T_{\mathrm{encode}}^{\mathrm{free}} \approx \frac{1}{3}\, T_{\mathrm{encode}}^{\mathrm{coded}},
\]
yielding a $\sim 3\times$ speedup.

\textbf{Decoding.} The standard pipeline
$\text{Entropy Decode} \to \text{Dequantize} \to \text{Inverse Transform}$
has its bottleneck in sequential entropy decoding. Fixed-length decoding
reads all indices in parallel:
\[
T_{\mathrm{decode}}^{\mathrm{free}} \approx \frac{1}{5}\, T_{\mathrm{decode}}^{\mathrm{coded}},
\]
yielding a $\sim 5\times$ speedup. The larger decode speedup reflects the
fact that entropy decoding is a more severe bottleneck on the decode path
(it must be perfectly sequential and bit-exact).

\textbf{Rate cost.} By Theorem~\ref{thm:rate-penalty}, the rate penalty is
$\Delta R = O(K^{-2/d})$, which is negligible for practical codebook sizes.
\end{corollary}

\begin{proof}
The speedup factors arise from the elimination of sequential dependencies:

\textbf{Encoding.} In the entropy-coded pipeline, encoding time is dominated by:
\begin{itemize}[nosep]
  \item Transform: $O(N)$ parallelizable (convolution/attention).
  \item Quantization: $O(N)$ parallelizable (nearest-neighbor lookup).
  \item Entropy coding: $O(N)$ \emph{sequential} (arithmetic coding processes
        symbols one at a time, maintaining a running probability context).
\end{itemize}
The entropy coding step typically accounts for $\sim 2/3$ of the total
encoding time (Cao et al., 2026, Table 3). Removing it reduces:
\[
T_{\mathrm{encode}}^{\mathrm{free}}
= T_{\mathrm{transform}} + T_{\mathrm{quant}}
\approx \frac{1}{3}\, T_{\mathrm{encode}}^{\mathrm{coded}}.
\]

\textbf{Decoding.} In the entropy-coded pipeline, decoding time is dominated by:
\begin{itemize}[nosep]
  \item Entropy decoding: $O(N)$ \emph{sequential}, $\sim 4/5$ of total.
  \item Dequantization + inverse transform: $O(N)$ parallelizable.
\end{itemize}
The entropy decoding step accounts for $\sim 4/5$ of decoding time because
arithmetic decoding is inherently sequential (each symbol depends on the
accumulated bitstream state). Removing it:
\[
T_{\mathrm{decode}}^{\mathrm{free}}
= T_{\mathrm{dequant}} + T_{\mathrm{inverse}}
\approx \frac{1}{5}\, T_{\mathrm{decode}}^{\mathrm{coded}}.
\]

The rate cost of this speedup is $\Delta R = O(K^{-2/d})$ bits/index
(Theorem~\ref{thm:rate-penalty}), which is negligible for practical
codebook sizes $K \geq 2^{10}$ and moderate dimensions $d \leq 64$.
\end{proof}

% =====================================================================
\section{Summary}
% =====================================================================

\begin{remark}[Summary of Results]
The three theorems establish a complete characterization:

\begin{itemize}[nosep]
  \item \textbf{Theorem 1}: Unconstrained VQ naturally produces near-uniform
        index distributions ($H \to \log_2 K$), making entropy coding
        unnecessary in the asymptotic regime.
  \item \textbf{Theorem 2}: The rate penalty of skipping entropy coding is
        exactly the entropy gap $\Delta R = \log_2 K - H(p)$, which vanishes
        as $O(K^{-2/d})$.
  \item \textbf{Theorem 3}: Entropy-free coding is optimal iff indices are
        uniform, independent, and context-irrelevant. Correlations can be
        removed via parallelizable context-conditioned transforms.
  \item \textbf{Corollary}: The practical payoff is $3\times$ encode and
        $5\times$ decode speedup at negligible rate cost.
\end{itemize}

These results formalize the EF-LIC insight: unconstrained VQ + fixed-length
coding achieves near-optimal compression with dramatically faster encoding
and decoding by eliminating the sequential entropy coding bottleneck.
\end{remark}

% =====================================================================
%  References
% =====================================================================
\begin{thebibliography}{9}

\bibitem{cao2026}
J.~Cao et al.,
\newblock ``EF-LIC: Entropy-Free Learned Image Compression,''
\newblock in \emph{Proc.\ ICML}, 2026.

\bibitem{gersho1992}
A.~Gersho and R.~M.~Gray,
\newblock \emph{Vector Quantization and Signal Compression},
\newblock Springer, 1992.

\bibitem{cover2006}
T.~M.~Cover and J.~A.~Thomas,
\newblock \emph{Elements of Information Theory},
\newblock 2nd ed., Wiley, 2006.

\bibitem{gray1998}
R.~M.~Gray and D.~L.~Neuhoff,
\newblock ``Quantization,''
\newblock \emph{IEEE Trans.\ Inf.\ Theory}, vol.~44, no.~6, pp.~2325--2383, 1998.

\bibitem{graf2000}
S.~Graf and H.~Luschgy,
\newblock \emph{Foundations of Quantization for Probability Distributions},
\newblock Springer, 2000.

\end{thebibliography}

\end{document}